% =============================================================
% เอกสารคู่มือการใช้งาน "ไดโน" (Dino)
% โครงสร้างอ้างอิงจากคู่มือการใช้งานของโครงงาน "ไปนำแหน่"
% (ปก > คำนำ > สารบัญ > สารบัญภาพ > บทที่ 1 บทนำ > บทที่ 2 คู่มือการใช้งาน > บทที่ 3 การติดตั้งระบบสำหรับพัฒนา)
% คอมไพล์ด้วย XeLaTeX:  cd manual && xelatex main.tex  (รัน 2 รอบ)
% =============================================================
\documentclass{cskkuproject}
\usepackage{needspace}
\usepackage{xcolor}

% ความกว้างป้าย "(1)" สำหรับรายการแบบ (1)(2)(3) — ใช้ค่าเดียวกับ report/main.tex
\newlength{\parenlabelwidth}
\settowidth{\parenlabelwidth}{(1)}

% รายการแบบ (1)(2)(3) ให้บรรทัดแรกเริ่มตรงกับย่อหน้าปัจจุบัน
\newlist{parenlist}{enumerate}{1}
\setlist[parenlist]{label=(\arabic*),leftmargin=0pt,
  itemindent=\dimexpr\parindent+\parenlabelwidth+1em\relax,
  align=right,labelwidth=\parenlabelwidth,labelsep=1em}

% รายการเลขลำดับธรรมดา 1. 2. 3. (ไม่ใช้รูปแบบ 2.1.1. ของคลาส)
\newlist{numlist}{enumerate}{1}
\setlist[numlist]{label=\arabic*.,leftmargin=*,labelindent=\dimexpr\parindent+1em\relax,
  labelwidth=2em,labelsep=0.5em,itemsep=0pt,topsep=0.3\baselineskip}

% -------------------------------------------------------------
% ข้อความที่ยังต้องยืนยัน/กรอกเพิ่ม แสดงเป็นตัวหนาสีแดงในวงเล็บเหลี่ยม
% ค้นหาในไฟล์ด้วยคำว่า \fillin เพื่อไล่ตรวจก่อนส่งเล่มจริง
% -------------------------------------------------------------
\newcommand{\fillin}[1]{{\bfseries\color{red}[#1]}}

% -------------------------------------------------------------
% ภาพหน้าจอ: \screenshot[ความกว้าง]{ชื่อไฟล์ไม่รวมนามสกุล}{คำอธิบายภาพ}
% วางไฟล์ภาพจริงที่ images/screens/<ชื่อไฟล์>.png ระบบจะแทนกรอบว่างให้อัตโนมัติ
% -------------------------------------------------------------
\newcommand{\screenshot}[3][0.85]{%
  \begin{figure}[H]
    \centering
    \setlength{\fboxsep}{0pt}%
    \IfFileExists{images/screens/#2.png}{%
      \fbox{\includegraphics[width=#1\textwidth]{images/screens/#2.png}}%
    }{\IfFileExists{images/screens/#2.jpg}{%
      \fbox{\includegraphics[width=#1\textwidth]{images/screens/#2.jpg}}%
    }{%
      \fbox{\parbox[c][5.5cm][c]{#1\textwidth}{\centering\fillin{ภาพหน้าจอ: images/screens/#2.png}}}%
    }}
    \caption{#3}
    \label{fig:#2}
  \end{figure}
}

% -------------------------------------------------------------
% ข้อมูลโครงงาน — ใช้ชื่อตัวแปรเดียวกับ report/main.tex
% -------------------------------------------------------------
\titleTH{ไดโน : เว็บแอพลิเคชันศูนย์กลางของการท่องเที่ยวจังหวัดขอนแก่น}
\titleEN{Dino: The central web application for tourism in Khon Kaen}

\firstStudentID{663380203-9}
\firstStudentPrefixTH{นาย}
\firstStudentPrefixEN{Mr.}
\firstStudentNameTH{คเณศ บัวสด}
\firstStudentNameEN{Khanet Buasod}

\secondStudentID{663380228-3}
\secondStudentPrefixTH{นาย}
\secondStudentPrefixEN{Mr.}
\secondStudentNameTH{ภูริ ชาริชา}
\secondStudentNameEN{Phuri Charicha}

\degreeTH{วิทยาศาสตรบัณฑิต}
\degreeEN{Bachelor of Science}
\majorTH{วิทยาการคอมพิวเตอร์}
\majorEN{Computer Science}

\advisorTH{อ.ธนพล ตั้งชูพงศ์}
\advisorEN{Thanapon Tangchupong}

\academicYearTH{2569}
\academicYearEN{2026}

\courseCode{CP354 771}
\courseName{โครงงานทางวิทยาการคอมพิวเตอร์ 1}
\semester{1}

\begin{document}

% ==========================================
% หน้าปก (ตามรูปแบบปกคู่มือการใช้งาน)
% ==========================================
\makeatletter
\begin{titlepage}
\thispagestyle{empty}
\centering
\vspace*{1cm}

\includegraphics[scale=0.5]{./theme/kku-logo.png}

\vspace{1.2cm}
{\Large\bfseries \@courseCode{} \@courseName\par}
\vspace{0.6cm}
{\Large\bfseries เอกสารคู่มือการใช้งาน\par}
\vspace{0.6cm}
{\large \@titleTH\par}
\vspace{0.3cm}
{\large \@titleEN\par}

\vspace{1.5cm}
{\large โดย\par}
\vspace{0.3cm}
{\large \@firstStudentID~\@firstStudentPrefixTH\@firstStudentNameTH\par}
{\large \@secondStudentID~\@secondStudentPrefixTH\@secondStudentNameTH\par}

\vspace{1cm}
{\large อาจารย์ที่ปรึกษา: \@advisorTH\par}

\vfill
{\normalsize รายงานนี้เป็นส่วนหนึ่งของการศึกษาวิชา \@courseCode{} \@courseName\par}
{\normalsize ภาคเรียนที่ \@semester{} ปีการศึกษา \@academicYearTH\par}
{\normalsize สาขาวิชา\@majorTH{} วิทยาลัยการคอมพิวเตอร์\par}
{\normalsize มหาวิทยาลัยขอนแก่น\par}
{\normalsize (เดือน ตุลาคม พ.ศ. \@academicYearTH)\par}
\end{titlepage}
\makeatother

% เลขหน้าส่วนต้นเป็นอักษรไทย (ก, ข, ค, ...)
\pagenumbering{arabic}
\setcounter{page}{0}
\renewcommand{\thepage}{\ThaiAlph{\value{page}}}

% ==========================================
% คำนำ
% ==========================================
% คำนำ — ปรับจากคำนำของคู่มือ "ไปนำแหน่" ให้ตรงกับฟังก์ชันของระบบไดโน
\chapter*{คำนำ}
\addcontentsline{toc}{chapter}{คำนำ}

เอกสารคู่มือการใช้งานเว็บแอปพลิเคชัน ``ไดโน'' ได้จัดทำขึ้นโดยมีวัตถุประสงค์เพื่อให้ผู้ใช้งานสามารถใช้งานระบบได้อย่างถูกต้องและเกิดประสิทธิภาพสูงสุด คู่มือนี้ประกอบด้วยขั้นตอนการเข้าใช้งาน การสมัครสมาชิกและเข้าสู่ระบบ การค้นหาสถานที่ท่องเที่ยวและกิจกรรม การสอบถามข้อมูลกับแชตบอต การสร้างและปรับแผนการเดินทาง การสะสมพอยท์ด้วยการสแกน QR Code ตลอดจนการใช้งานในส่วนของผู้ดูแลระบบ และขั้นตอนการติดตั้งระบบสำหรับนักพัฒนา ผู้จัดทำหวังเป็นอย่างยิ่งว่าคู่มือฉบับนี้จะเป็นประโยชน์ต่อผู้ใช้งานทุกกลุ่ม

\vspace{1.5cm}
\begin{flushright}
ผู้จัดทำ\\[0.5em]
นาย คเณศ บัวสด\\[0.5em]
นาย ภูริ ชาริชา
\end{flushright}
\clearpage


% ==========================================
% สารบัญ / สารบัญภาพ / สารบัญตาราง
% ==========================================
\tableofcontents
\clearpage
\listoffigures
\clearpage
\listoftables
\clearpage

% ==========================================
% เนื้อหา 3 บท
% ==========================================
\pagenumbering{arabic}
\setcounter{page}{1}
\renewcommand{\thepage}{\arabic{page}}

\pagestyle{fancy}
\fancyhf{}
\fancyhead[R]{\thepage}
\setlength{\headheight}{15pt}

\chapter{บทนำ}
% =============================================================
% บทที่ 1 บทนำ — เครื่องมือที่ใช้พัฒนา และขอบเขต/ข้อจำกัดของระบบ
% ข้อมูลตรวจจาก main repo (README.md, CLAUDE.md, package.json) และเว็บจริง dinokhonkaen.app
% =============================================================

\section{เครื่องมือที่ใช้พัฒนา}

เว็บแอปพลิเคชัน ``ไดโน'' (Dino Khon Kaen) เผยแพร่ใช้งานจริงที่ https://dinokhonkaen.app พัฒนาเป็นบริการอิสระ 3 ส่วน ได้แก่ ส่วนติดต่อผู้ใช้งาน บริการข้อมูลหลัก และ AI Service ซึ่งแยกโฟลเดอร์และทำงานแยกจากกัน โดยใช้เครื่องมือดังนี้

\subsection{React และ Vite}

ใช้พัฒนาส่วนติดต่อผู้ใช้งาน (โฟลเดอร์ \texttt{frontend}) ซึ่งเป็นเว็บแอปพลิเคชันชุดเดียวที่ให้บริการทั้งฝั่งผู้ใช้งานและหน้าผู้ดูแลระบบ ใช้ React 18 ร่วมกับ react-router-dom และใช้ Vite เป็นเครื่องมือพัฒนาและสร้างไฟล์สำหรับใช้งานจริง

\subsection{Node.js และ Express}

ใช้พัฒนาบริการข้อมูลหลัก (โฟลเดอร์ \texttt{backend}) ให้บริการ REST API ภายใต้เส้นทาง \texttt{/api} สำหรับจัดการข้อมูลสถานที่ กิจกรรม ฐานความรู้ QR Code ของรางวัล พอยท์ ทริป และระบบสมาชิก พร้อมทั้งเชื่อมต่อฐานข้อมูลผ่านไลบรารี knex

\subsection{Python และ FastAPI}

ใช้พัฒนา AI Service (โฟลเดอร์ \texttt{chatbot-service}) ซึ่งรับผิดชอบแชตบอต ระบบวางแผนทริป การสกัดข้อมูลกิจกรรมจากข้อความโพสต์ และการสร้าง Embedding โดยเรียกโมเดลภาษาขนาดใหญ่ผ่าน KKU AI Gateway

\subsection{PostgreSQL บน Supabase}

ใช้เป็นฐานข้อมูลหลักของระบบ พร้อมส่วนขยาย pgvector สำหรับการค้นคืนข้อมูลเชิงเวกเตอร์ โดยโครงสร้างตารางอยู่ในไฟล์ migration ของโฟลเดอร์ \texttt{backend/supabase/migrations}

\subsection{บริการภายนอก}

ระบบใช้ Cloudinary สำหรับเก็บรูปภาพ Google Places API สำหรับนำเข้าข้อมูลสถานที่ Google Maps JavaScript API สำหรับแสดงแผนที่ และ Gmail SMTP สำหรับส่งอีเมลยืนยันและรีเซ็ตรหัสผ่าน

\subsection{Visual Studio Code}

ใช้เป็นเครื่องมือสำหรับเขียนและแก้ไขซอร์สโค้ดของโปรแกรม

\section{ส่วนประมวลผลปัญญาประดิษฐ์ (AI Service)}

\subsection{ภาพรวมและหน้าที่}

AI Service คือบริการที่พัฒนาด้วย Python และ FastAPI ในโฟลเดอร์ \texttt{chatbot-service} ทำหน้าที่ประมวลผลทุกฟังก์ชันที่ใช้ปัญญาประดิษฐ์ของระบบ โดยเรียกใช้โมเดลภาษาขนาดใหญ่ (Large Language Model: LLM) ผ่าน KKU AI Gateway ซึ่งมีรูปแบบการเรียกใช้เข้ากันได้กับ OpenAI และสามารถเปลี่ยนไปใช้ Vertex AI ของ Google แทนได้ด้วยการตั้งค่า ฟังก์ชันที่ AI Service ให้บริการแสดงดังตารางที่ \ref{tab:ai-functions}

\begin{table}[H]
\caption{ฟังก์ชันของ AI Service}
\label{tab:ai-functions}
\centering
\small
\begin{tabular}{|L{3.0cm}|L{3.3cm}|L{3.0cm}|L{3.1cm}|}
\hline
\textbf{ฟังก์ชัน} & \textbf{เส้นทาง} & \textbf{ผู้เรียกใช้} & \textbf{ใช้ LLM} \\
\hline
แชตบอตน้องไดโน & \texttt{/chat/} & เบราว์เซอร์ของผู้ใช้ & ใช้ (ปรับคำถาม และตอบ) \\
\hline
วางแผนทริป & \texttt{/trip/llm} & เบราว์เซอร์ของผู้ใช้ & ใช้ (จัดแผน และตรวจคุณภาพ) \\
\hline
สกัดข้อมูลกิจกรรมจากโพสต์ & \texttt{/events/extract} & backend (ผู้ดูแลระบบ) & ใช้ \\
\hline
สร้างคำอธิบายสถานที่ & \texttt{/admin/places/describe} & backend (ผู้ดูแลระบบ) & ใช้ \\
\hline
สร้างดัชนีค้นหา (Embedding) & \texttt{/admin/reindex} & backend (ผู้ดูแลระบบ) & ไม่ใช้ \\
\hline
\end{tabular}
\end{table}

เส้นทางสำหรับผู้ดูแลระบบของ AI Service ไม่มีการตรวจสิทธิ์ในตัวเอง จึงถูกกำหนดให้เรียกผ่าน backend ซึ่งเป็นผู้ตรวจสิทธิ์ผู้ดูแลระบบเท่านั้น และไม่เปิดให้เบราว์เซอร์เรียกโดยตรง

\subsection{ข้อมูลและโมเดลที่ใช้}

AI Service ใช้ข้อมูล 3 ชุดในฐานข้อมูล ได้แก่ ข้อมูลสถานที่ ฐานความรู้ของแชตบอต และข้อมูลกิจกรรม โดยแปลงข้อความของแต่ละรายการเป็นเวกเตอร์ (Embedding) ขนาด 384 มิติ ด้วยโมเดล \texttt{paraphrase-multilingual-MiniLM-L12-v2} ซึ่งรองรับภาษาไทยและทำงานบนเครื่องของ AI Service เอง ไม่ต้องเรียกบริการภายนอก จากนั้นเก็บไว้ในคอลัมน์ \texttt{embedding} ของฐานข้อมูล PostgreSQL ที่เปิดใช้ส่วนขยาย pgvector เพื่อใช้ค้นหาตามความหมาย

การค้นหาใช้ 2 วิธีร่วมกัน (hybrid search) คือ ค้นตามความคล้ายของเวกเตอร์ และค้นตามคำสำคัญในชื่อและป้ายกำกับของสถานที่ เพื่อไม่ให้พลาดรายการที่ตอบคำถามได้จริงแต่เวกเตอร์ไม่ใกล้เคียงมากพอ นอกจากนี้ระบบขยายคำถามด้วยคำพ้องความหมายก่อนค้นหา ส่วนกิจกรรมที่ถูกยกเลิกหรือสิ้นสุดแล้วจะไม่ถูกนำมาตอบ

เมื่อผู้ดูแลระบบเพิ่มหรือแก้ไขข้อมูล รายการนั้นจะถูกทำเครื่องหมายว่ารอสร้างเวกเตอร์ใหม่ และต้องกดปุ่ม ``อัปเดตดัชนีค้นหา'' ในหน้าแดชบอร์ดผู้ดูแลระบบ เพื่อให้ระบบสร้างเวกเตอร์เฉพาะรายการที่รออยู่ การสร้างเวกเตอร์ใหม่ทั้งหมดใช้สคริปต์ \texttt{scripts/embed\_content.py} และต้องทำเมื่อเปลี่ยนโมเดล Embedding

\subsection{การทำงานของแชตบอตน้องไดโน}

แชตบอตใช้เทคนิค Retrieval-Augmented Generation (RAG) คือ ค้นข้อมูลที่เกี่ยวข้องจากฐานข้อมูลก่อน แล้วให้ LLM เรียบเรียงคำตอบโดยอ้างอิงเฉพาะข้อมูลที่ค้นได้ ขั้นตอนการทำงานในแต่ละคำถามมีดังนี้

\begin{numlist}
\item รับข้อความของผู้ใช้ (ไม่เกิน 500 ตัวอักษร) พร้อมประวัติสนทนาล่าสุดไม่เกิน 4 ข้อความ ซึ่งฝั่งผู้ใช้เป็นผู้เก็บประวัติ ส่วน AI Service ไม่เก็บสถานะการสนทนา
\item หากมีประวัติสนทนา ระบบให้ LLM เรียบเรียงคำถามใหม่ให้สมบูรณ์ในตัวเอง เช่น คำถามต่อเนื่องว่า ``แล้วร้านที่สองล่ะ'' จะถูกขยายเป็นคำถามที่ระบุชื่อสถานที่ชัดเจน เพื่อให้ค้นหาได้ถูกต้อง
\item ค้นข้อมูลที่เกี่ยวข้องอย่างละไม่เกิน 3 รายการ จากสถานที่ ฐานความรู้ และกิจกรรม
\item ส่งข้อมูลที่ค้นได้และประวัติสนทนาให้ LLM ตอบ โดยกำหนดให้ตอบจากข้อมูลที่ได้รับเท่านั้น และให้ระบุแหล่งที่มาของคำตอบกำกับไว้ในรูปแบบแท็ก
\item ระบบอ่านแท็กดังกล่าวเพื่อตัดสินว่าคำตอบอ้างอิงข้อมูลชุดใด แล้วแนบการ์ดสถานที่หรือกิจกรรมที่ใช้อ้างอิงไปกับคำตอบ หากไม่พบข้อมูลที่เกี่ยวข้อง แชตบอตจะตอบข้อความแจ้งว่าไม่มีข้อมูลในส่วนนี้ แทนการคาดเดา
\item ส่งคำตอบกลับแบบสตรีม (Server-Sent Events) เป็นช่วง ๆ เพื่อให้ผู้ใช้เห็นคำตอบทยอยแสดงโดยไม่ต้องรอจนเสร็จ
\end{numlist}

\subsection{การทำงานของระบบวางแผนทริป}

ระบบวางแผนทริปเป็นแบบผสม (hybrid) คือ ให้ LLM ทำหน้าที่เลือกและจัดลำดับสถานที่ ส่วนการคำนวณระยะทาง เวลาเดินทาง เวลาเปิด-ปิด และค่าใช้จ่ายใช้โปรแกรมที่คำนวณได้แน่นอน (deterministic) เพื่อไม่ให้ผลลัพธ์ที่เป็นตัวเลขผิดพลาดจากการสร้างข้อความของโมเดล ขั้นตอนการทำงานมีดังนี้

\begin{numlist}
\item รับข้อมูลจากแบบฟอร์ม 4 ขั้นตอน ได้แก่ จำนวนวันและวันเริ่มต้น ที่พัก สถานที่ที่ต้องไป ความสนใจ งบประมาณ จังหวะการเที่ยว ขอบเขตพื้นที่ และช่วงเวลาของแต่ละวัน
\item สร้างรายการสถานที่ตัวเลือก โดยค้นตามความสนใจแต่ละข้อ กรองตามระดับงบประมาณและขอบเขตพื้นที่ และเผื่อจำนวนตัวเลือกให้มากกว่าที่ต้องใช้ เพื่อให้ระบบมีสถานที่สำรองเติมช่วงเวลาที่ว่าง
\item ให้ LLM เลือกและจัดสถานที่เป็นรายวัน โดยกำหนดให้ตอบเป็น JSON และมีข้อจำกัด เช่น ร้านอาหารไม่เกิน 2 แห่งต่อวัน และไม่นำคาเฟ่มาเรียงต่อกันหลายแห่ง
\item โปรแกรมคำนวณตารางเวลาจริง ได้แก่ ลำดับการเดินทาง เวลาถึงและเวลาออก ระยะทางรวม และค่าใช้จ่ายโดยประมาณ พร้อมตรวจเวลาเปิด-ปิดของสถานที่ และเติมสถานที่สำรองหากวันใดยังว่างอยู่
\item ให้ LLM อีกชุดทำหน้าที่ผู้ตรวจ (judge) ประเมินว่าแผนเหมาะกับจังหวะการเที่ยวและตรงกับความสนใจของผู้ใช้หรือไม่ หากไม่ผ่าน ระบบจะสร้างแผนใหม่พร้อมข้อเสนอแนะ โดยทำซ้ำได้ไม่เกิน 2 รอบ
\item ส่งผลลัพธ์กลับพร้อมสรุปเหตุผลของการจัดแผน และแจ้งข้อควรระวัง เช่น สถานที่ที่ระบุแต่จัดเข้าแผนไม่ได้ หรือช่วงเดินทางที่ตรงกับวันหยุดนักขัตฤกษ์
\end{numlist}

หากโมเดลตอบรูปแบบไม่ถูกต้องหรือการวางแผนล้มเหลว ระบบจะส่งข้อความ ``วางแผนทริปไม่สำเร็จ กรุณาลองใหม่อีกครั้ง'' กลับไป โดยไม่สร้างแผนสำรองด้วยการคาดเดา

\subsection{การสกัดข้อมูลกิจกรรมและการสร้างคำอธิบายสถานที่}

ฟังก์ชันสำหรับผู้ดูแลระบบ 2 ฟังก์ชันใช้ LLM ช่วยลดงานกรอกข้อมูล โดยผลลัพธ์ทุกครั้งต้องให้ผู้ดูแลระบบตรวจสอบก่อนบันทึก

\begin{parenlist}
\item การสกัดข้อมูลกิจกรรม รับข้อความโพสต์ภาษาไทยแล้วให้ LLM แยกเป็นช่องข้อมูลของแบบฟอร์มกิจกรรม เช่น ชื่องาน ช่วงวันที่ สถานที่ ค่าเข้าชม และผู้จัดงาน โดยกำหนดให้ใช้เฉพาะข้อมูลที่ปรากฏในโพสต์ หากไม่ระบุให้เว้นว่าง และระบบตรวจรูปแบบวันที่ซ้ำ ตัดช่วงวันที่ที่ผิดพลาดออก
\item การสร้างคำอธิบายสถานที่ นำชื่อ หมวดหมู่ ป้ายกำกับ และรีวิวจาก Google มาสรุปเป็นคำอธิบาย และหากกำหนดคีย์ Serper ไว้ จะใช้ผลค้นหาจากเว็บประกอบด้วย หากข้อมูลไม่พอระบบใช้ข้อความแม่แบบแทน และสามารถสลับไปใช้โมเดลสำรองเมื่อโควตารายวันของโมเดลหลักหมด
\end{parenlist}

\subsection{การจำกัดการใช้งานและการติดตามค่าใช้จ่าย}

เนื่องจากการเรียกใช้ LLM มีค่าใช้จ่ายตามปริมาณโทเคน AI Service จึงมีมาตรการควบคุม ดังนี้

\begin{parenlist}
\item จำกัดจำนวนคำขอต่อผู้ใช้ (ตาม IP) โดยค่าเริ่มต้นคือ แชตบอต 20 ครั้งต่อนาทีและ 200 ครั้งต่อชั่วโมง และการวางแผนทริป 5 ครั้งต่อ 10 นาทีและ 20 ครั้งต่อชั่วโมง
\item จำกัดจำนวนคำขอรวมทุกผู้ใช้ ค่าเริ่มต้นคือ แชตบอต 120 ครั้งต่อนาทีและ 3,000 ครั้งต่อชั่วโมง และการวางแผนทริป 30 ครั้งต่อ 10 นาทีและ 200 ครั้งต่อชั่วโมง เพื่อกำหนดเพดานค่าใช้จ่ายสูงสุด เมื่อเกินจะตอบกลับสถานะ 429
\item จำกัดความยาวข้อความที่รับเข้า เช่น ข้อความแชตไม่เกิน 500 ตัวอักษร และจำนวนวันของทริปไม่เกิน 14 วัน
\item บันทึกจำนวนโทเคนและค่าใช้จ่ายโดยประมาณของทุกคำขอลงไฟล์ \texttt{logs/usage.jsonl} และสรุปผลได้ด้วยสคริปต์ \texttt{scripts/usage\_report.py}
\end{parenlist}

ค่าจำกัดข้างต้นปรับได้ด้วยตัวแปรสภาพแวดล้อมชุด \texttt{RATE\_LIMIT\_*} ในไฟล์ \texttt{chatbot.env} ตามที่ระบุไว้ในไฟล์ \texttt{.env.example}

\section{ขอบเขตและข้อจำกัด}

\subsection{ขอบเขตของผู้ใช้งานทั่วไป (Guest)}

\begin{parenlist}
\item เข้าชมหน้าเว็บสาธารณะ ค้นหาและดูรายละเอียดสถานที่ท่องเที่ยว ร้านอาหาร คาเฟ่ และกิจกรรมได้
\item สอบถามข้อมูลการท่องเที่ยวกับน้องไดโน (แชตบอต) ได้
\item สร้างแผนการเดินทางด้วย AI ได้ แต่แผนจะสูญหายเมื่อปิดหน้าเว็บ เนื่องจากไม่ได้บันทึกลงบัญชี
\item ดูรายการสถานที่ร่วมรายการและของรางวัลในหน้าพอยท์สะสมได้ แต่สะสมพอยท์ไม่ได้
\end{parenlist}

\subsection{ขอบเขตของสมาชิก (Tourist)}

\begin{parenlist}
\item สมัครสมาชิกด้วยอีเมล ยืนยันอีเมล และเข้าสู่ระบบ
\item ใช้งานได้ทุกฟังก์ชันของผู้ใช้งานทั่วไป และบันทึกแผนทริปไว้ในเมนู ``ทริปของฉัน''
\item สแกน QR Code ณ สถานที่ร่วมรายการเพื่อสะสมพอยท์ และดูประวัติพอยท์
\item แลกของรางวัลที่เคาน์เตอร์โดยแจ้งเจ้าหน้าที่
\item แก้ไขข้อมูลโปรไฟล์ เปลี่ยนรหัสผ่านหรืออีเมล ส่งกิจกรรมใหม่เข้าระบบ และแจ้งข้อมูลสถานที่หรือกิจกรรมที่ไม่ถูกต้อง
\end{parenlist}

\subsection{ขอบเขตของผู้ดูแลระบบ (Admin)}

\begin{parenlist}
\item เข้าสู่ระบบผ่านหน้าเข้าสู่ระบบสำหรับผู้ดูแลระบบ
\item ดูภาพรวมของระบบผ่านแดชบอร์ด
\item เพิ่ม แก้ไข และลบข้อมูลสถานที่ กิจกรรม และฐานความรู้ของแชตบอต รวมถึงปรับปรุงข้อมูลสถานที่ให้ตรงกับ Google Places
\item ตรวจสอบรายงานข้อมูลที่ไม่ถูกต้องและคำขอกิจกรรมจากผู้ใช้
\item จัดการ QR Code ของรางวัล และดำเนินการแลกของรางวัลที่เคาน์เตอร์
\item จัดการผู้ใช้ ดูทริปที่ผู้ใช้สร้าง และดูผลความพึงพอใจ
\end{parenlist}

\subsection{ข้อจำกัดของระบบ}

\begin{parenlist}
\item ข้อมูลสถานที่ครอบคลุมจังหวัดขอนแก่น โดยข้อมูลหลักมาจาก Google Places API และข้อมูลที่ผู้ดูแลระบบเพิ่มเอง จึงอาจไม่เป็นปัจจุบันหากยังไม่ได้ปรับปรุงข้อมูล
\item แชตบอตและระบบวางแผนทริปใช้โมเดลภาษาขนาดใหญ่ แม้จำกัดการตอบให้อ้างอิงฐานความรู้ แต่ยังมีโอกาสตอบคลาดเคลื่อน และอาจใช้งานไม่ได้ชั่วคราวเมื่อโควตาของ KKU AI Gateway หมด
\item QR Code แต่ละจุดให้พอยท์ได้ครั้งเดียวต่อบัญชี และการแลกของรางวัลต้องทำที่เคาน์เตอร์โดยเจ้าหน้าที่ ระบบยังไม่รองรับการแลกของรางวัลด้วยตนเอง
\item ผู้ดูแลระบบเป็นบัญชีที่สร้างโดยนักพัฒนาผ่านสคริปต์ ไม่สามารถสมัครผ่านหน้าเว็บไซต์ได้
\item ระบบเป็นเว็บแอปพลิเคชันที่ใช้ผ่านเบราว์เซอร์ ไม่มีแอปพลิเคชันบนมือถือแยกต่างหาก และการสแกน QR Code ต้องอนุญาตให้เว็บเข้าถึงกล้อง
\end{parenlist}


\chapter{คู่มือการใช้งาน}
% =============================================================
% บทที่ 2 คู่มือการใช้งาน
% เนื้อหาตรวจจากเว็บจริง https://dinokhonkaen.app และซอร์สโค้ดใน main repo/frontend
% ภาพหน้าจอ: images/screens/<ชื่อ>.jpg|png  (ถ้ายังไม่มีไฟล์ จะแสดงกรอบแดงแทน)
% =============================================================

\section{คู่มือการใช้งาน}

``ไดโน'' เว็บแอปพลิเคชันศูนย์กลางของการท่องเที่ยวจังหวัดขอนแก่น เข้าใช้งานได้ที่ https://dinokhonkaen.app แบ่งผู้ใช้งานออกเป็น 3 กลุ่ม ได้แก่

\begin{numlist}
\item ผู้ใช้งานทั่วไป (Guest)
\item สมาชิก (Tourist)
\item ผู้ดูแลระบบ (Admin)
\end{numlist}

% -------------------------------------------------------------
\subsection{เข้าใช้งานกรณีเป็นผู้ใช้งานทั่วไป}
% -------------------------------------------------------------

แถบเมนูผู้ใช้งานทั่วไป อยู่ด้านบนของทุกหน้า ดังภาพที่ \ref{fig:guest-home}

\screenshot{guest-home}{หน้าแรกและแถบเมนูผู้ใช้งานทั่วไป}

จากภาพที่ \ref{fig:guest-home} รายละเอียดแถบเมนูมีดังนี้

\begin{numlist}
\item หน้าแรก แสดงช่องค้นหา กิจกรรมและเทศกาลที่กำลังจะมาถึง และสถานที่แนะนำ
\item สถานที่ท่องเที่ยว แสดงรายการสถานที่ทั้งหมดพร้อมตัวกรองตามหมวดหมู่
\item กิจกรรม \& เทศกาล แสดงรายการกิจกรรมและเทศกาลในขอนแก่น
\item วางแผนทริป AI ใช้สร้างแผนการเดินทางให้อัตโนมัติ
\item พอยท์สะสม แสดงสถานที่ร่วมรายการและของรางวัล
\item เข้าสู่ระบบ และสมัครสมาชิก อยู่มุมขวาบน
\end{numlist}

เมื่อเปิดเว็บไซต์ครั้งแรก ระบบอาจแสดงหน้าต่างต้อนรับจากน้องไดโน สามารถกดปุ่ม ``เริ่มสำรวจขอนแก่น'' หรือ ``ไม่แสดงอีก 7 วัน'' เพื่อปิดได้ และมีปุ่มรูปน้องไดโนมุมขวาล่างของทุกหน้าสำหรับเปิดแชตบอต

\subsubsection{เมนูสมัครสมาชิก กรณีเป็นผู้ใช้งานทั่วไป}

การสมัครสมาชิกทำได้โดยคลิกปุ่ม ``สมัครสมาชิก'' มุมขวาบนของหน้าเว็บ ระบบจะแสดงแบบฟอร์มสมัครสมาชิกดังภาพที่ \ref{fig:guest-signup-form}

\screenshot{guest-signup-form}{หน้าฟอร์มสมัครสมาชิก}

จากภาพที่ \ref{fig:guest-signup-form} ผู้ใช้กรอกข้อมูล ได้แก่ คำนำหน้า ชื่อ นามสกุล เพศ วันเกิด อาชีพ อีเมล เบอร์โทรศัพท์ รหัสผ่านและยืนยันรหัสผ่าน และที่อยู่ปัจจุบัน (จังหวัด อำเภอ ตำบล) โดยรหัสผ่านต้องมีความยาวอย่างน้อย 8 ตัวอักษร ประกอบด้วยตัวพิมพ์ใหญ่ ตัวพิมพ์เล็ก และตัวเลข ระบบแสดงผลการตรวจสอบเงื่อนไขทีละข้อใต้ช่องรหัสผ่าน จากนั้นต้องทำเครื่องหมายยินยอมให้เก็บและใช้ข้อมูลส่วนบุคคลตามนโยบายความเป็นส่วนตัว แล้วกดปุ่ม ``สมัครสมาชิก''

หลังสมัครสำเร็จ ระบบจะส่งอีเมลยืนยันไปยังที่อยู่อีเมลที่กรอก ผู้ใช้ต้องคลิกลิงก์ในอีเมลเพื่อยืนยันบัญชี \fillin{ยืนยันกับระบบจริงว่าต้องยืนยันอีเมลก่อนเข้าสู่ระบบหรือไม่}

\subsubsection{เมนูเข้าสู่ระบบ กรณีเป็นผู้ใช้งานทั่วไป}

ผู้ใช้ที่มีบัญชีแล้วคลิกเมนู ``เข้าสู่ระบบ'' มุมขวาบน ระบบจะแสดงแบบฟอร์มเข้าสู่ระบบดังภาพที่ \ref{fig:guest-login-form}

\screenshot{guest-login-form}{หน้าฟอร์มเข้าสู่ระบบ}

ผู้ใช้กรอกอีเมลและรหัสผ่านที่ลงทะเบียนไว้ แล้วกดปุ่ม ``เข้าสู่ระบบ'' หากยังไม่มีบัญชีสามารถคลิก ``สมัครสมาชิก'' ด้านล่างปุ่ม เพื่อไปยังหน้าสมัครสมาชิกได้ทันที

หากลืมรหัสผ่าน ให้คลิก ``ลืมรหัสผ่าน?'' เพื่อเข้าสู่หน้าดังภาพที่ \ref{fig:forgot-password} กรอกอีเมลที่ใช้สมัครแล้วกด ``ส่งลิงก์รีเซ็ตรหัสผ่าน'' ระบบจะส่งลิงก์สำหรับตั้งรหัสผ่านใหม่ไปยังอีเมลนั้น

\screenshot{forgot-password}{หน้าลืมรหัสผ่าน}

\subsubsection{เมนูสถานที่ท่องเที่ยว กรณีเป็นผู้ใช้งานทั่วไป}

ผู้ใช้เรียกดูสถานที่ได้จากเมนู ``สถานที่ท่องเที่ยว'' ระบบจะแสดงรายการสถานที่ในรูปแบบการ์ด พร้อมช่องค้นหาและตัวกรองหมวดหมู่ ได้แก่ ทั้งหมด คาเฟ่ วัด พิพิธภัณฑ์ สวนสาธารณะ อุทยานแห่งชาติ ตลาด สถานที่ท่องเที่ยว และร้านอาหาร ดังภาพที่ \ref{fig:places-list}

\screenshot{places-list}{หน้ารายการสถานที่ท่องเที่ยว}

เมื่อคลิกการ์ดสถานที่ ระบบจะแสดงหน้ารายละเอียดดังภาพที่ \ref{fig:place-detail} ซึ่งประกอบด้วยหมวดหมู่ คะแนนและช่วงราคา แกลเลอรีรูปภาพ สรุปโดย Dino AI สิ่งอำนวยความสะดวก และรีวิวจากผู้เยี่ยมชม พร้อมปุ่ม ``เปิดเส้นทาง Google Maps'' และปุ่ม ``โทรติดต่อ''

\screenshot{place-detail}{หน้ารายละเอียดสถานที่ท่องเที่ยว}

\subsubsection{เมนูกิจกรรมและเทศกาล กรณีเป็นผู้ใช้งานทั่วไป}

เมนู ``กิจกรรม \& เทศกาล'' แสดงรายการกิจกรรมและเทศกาลในขอนแก่น พร้อมช่องค้นหาและตัวกรองสถานะ ได้แก่ ทั้งหมด กำลังจัดอยู่ เร็ว ๆ นี้ จบแล้ว และยกเลิก ดังภาพที่ \ref{fig:events-list}

\screenshot{events-list}{หน้ารายการกิจกรรมและเทศกาล}

เมื่อคลิกกิจกรรม ระบบจะแสดงรายละเอียด ได้แก่ รูปภาพ คำอธิบาย วันจัดงาน สถานที่จัดงาน ค่าเข้าชม ผู้จัดงาน กลุ่มที่เหมาะสม และแผนที่ ดังภาพที่ \ref{fig:event-detail} หากพบข้อมูลไม่ถูกต้อง สามารถกดปุ่ม ``แจ้งข้อมูลไม่ถูกต้อง'' ใต้รูปภาพเพื่อส่งให้ผู้ดูแลระบบตรวจสอบได้

\screenshot{event-detail}{หน้ารายละเอียดกิจกรรม}

\subsubsection{เมนูแชตบอต กรณีเป็นผู้ใช้งานทั่วไป}

ผู้ใช้คลิกปุ่มรูปน้องไดโนมุมขวาล่างเพื่อเปิดหน้าต่างแชตบอต พิมพ์คำถามเกี่ยวกับสถานที่ อาหาร หรือเทศกาลในช่องด้านล่าง แล้วกดปุ่มส่ง ดังภาพที่ \ref{fig:chatbot}

\screenshot[0.9]{chatbot}{หน้าต่างแชตบอตน้องไดโน}

ระบบจะค้นหาข้อมูลที่เกี่ยวข้องจากฐานความรู้และตอบกลับแบบทยอยปรากฏทีละข้อความ หากไม่พบข้อมูลที่เกี่ยวข้อง แชตบอตจะแจ้งว่าไม่มีข้อมูลแทนการตอบโดยไม่มีข้อมูลรองรับ ในกรณีที่ปรากฏข้อความ ``ระบบแชทขัดข้องชั่วคราว'' ให้ลองใหม่ภายหลัง

\subsubsection{เมนูวางแผนทริป AI กรณีเป็นผู้ใช้งานทั่วไป}

ผู้ใช้เลือกเมนู ``วางแผนทริป AI'' เพื่อสร้างแผนการเดินทางผ่านแบบฟอร์ม 4 ขั้นตอน โดยมีแถบแสดงความคืบหน้าด้านบน ขั้นตอนที่ 1 วันเดินทาง ดังภาพที่ \ref{fig:trip-step1}

\screenshot{trip-step1}{หน้าแบบฟอร์มสร้างแผนการเดินทาง ขั้นตอนที่ 1 วันเดินทาง}

จากภาพที่ \ref{fig:trip-step1} ผู้ใช้เลือกช่วงเวลาเดินทางจากปุ่มลัด ได้แก่ วันนี้ (เดย์ทริป) 2 วัน 3 วัน หรือสุดสัปดาห์นี้ หรือกำหนดวันเริ่มต้นและวันสิ้นสุดเอง ระบุที่พัก โดยพิมพ์ค้นหาหรือกด ``ตำแหน่งปัจจุบัน'' และระบุสถานที่ที่ต้องการไปแน่นอน (ถ้ามี) หากไม่ระบุที่พัก ระบบจะใช้ใจกลางเมืองขอนแก่นแทน จากนั้นกด ``ถัดไป''

ขั้นตอนที่ 2 ความสนใจ ผู้ใช้เลือกได้หลายข้อ ได้แก่ ธรรมชาติ วัฒนธรรม/ศาสนา ไดโนเสาร์ อาหารพื้นถิ่น คาเฟ่ ช้อปปิ้ง/หัตถกรรม และครอบครัว ดังภาพที่ \ref{fig:trip-step2}

\screenshot{trip-step2}{หน้าแบบฟอร์มสร้างแผนการเดินทาง ขั้นตอนที่ 2 ความสนใจ}

ขั้นตอนที่ 3 งบประมาณ ผู้ใช้เลือกงบประมาณโดยรวม (ประหยัด ปานกลาง หรูหรา) จังหวะการเที่ยว (สายชิลล์ กำลังดี หรือสายลุย) และขอบเขตพื้นที่ (เมือง หรือทั่วขอนแก่นซึ่งรวมสถานที่รอบนอก เช่น ภูเวียง อุบลรัตน์) ดังภาพที่ \ref{fig:trip-step3}

\screenshot{trip-step3}{หน้าแบบฟอร์มสร้างแผนการเดินทาง ขั้นตอนที่ 3 งบประมาณ}

ขั้นตอนที่ 4 ช่วงเวลา ผู้ใช้กำหนดเวลาเริ่มและเวลาสิ้นสุดของแต่ละวัน (ค่าเริ่มต้น 09:00 ถึง 18:00) แล้วกดปุ่ม ``สร้างแผนการเดินทาง'' ดังภาพที่ \ref{fig:trip-step4}

\screenshot{trip-step4}{หน้าแบบฟอร์มสร้างแผนการเดินทาง ขั้นตอนที่ 4 ช่วงเวลา}

เมื่อกดสร้างแผน ระบบแสดงหน้าสถานะกำลังประมวลผล ดังภาพที่ \ref{fig:trip-loading} จากนั้นแสดงแผนการเดินทางรายวันพร้อมแผนที่เส้นทาง ดังภาพที่ \ref{fig:trip-result}

\screenshot[0.7]{trip-loading}{หน้าสถานะกำลังประมวลผลแผนการเดินทาง}

\screenshot{trip-result}{หน้าแสดงผลแผนการเดินทาง}

ในหน้าผลลัพธ์ ผู้ใช้สามารถพิมพ์สิ่งที่อยากปรับ แล้วกดปุ่ม ``ปรับแผนการเดินทางอีกครั้ง'' เพื่อให้ระบบสร้างแผนใหม่ ทั้งนี้ ผู้ใช้งานทั่วไปที่ไม่ได้เข้าสู่ระบบจะไม่สามารถบันทึกแผนได้ แผนจะหายเมื่อปิดหน้า จึงควรเข้าสู่ระบบก่อนสร้างแผนเพื่อบันทึกไว้ในเมนู ``ทริปของฉัน''

\subsubsection{เมนูพอยท์สะสม กรณีเป็นผู้ใช้งานทั่วไป}

เมนู ``พอยท์สะสม'' สำหรับผู้ที่ยังไม่เข้าสู่ระบบ แสดงบัตรสะสมพอยท์ที่ระบุสถานที่ร่วมรายการพร้อมจำนวนพอยท์ที่ได้รับจากแต่ละที่ ขั้นตอนการสะสม 3 ขั้นตอน และรายการของรางวัลที่แลกได้ ดังภาพที่ \ref{fig:points-guest}

\screenshot{points-guest}{หน้าพอยท์สะสม กรณีผู้ใช้งานทั่วไป}

ผู้ใช้สามารถกด ``สมัครสมาชิกฟรี'' หรือ ``เข้าสู่ระบบ'' เพื่อเริ่มสะสมพอยท์ได้

% -------------------------------------------------------------
\subsection{เข้าใช้งานกรณีเข้าสู่ระบบของสมาชิก}
% -------------------------------------------------------------

เมื่อเข้าสู่ระบบแล้ว แถบเมนูจะเพิ่มเมนู ``ทริปของฉัน'' และแทนที่ปุ่มเข้าสู่ระบบด้วยเมนูผู้ใช้มุมขวาบน ดังภาพที่ \ref{fig:member-navbar}

\screenshot{member-navbar}{แถบเมนูสมาชิก}

\subsubsection{เมนูพอยท์สะสม กรณีเป็นสมาชิก}

หน้า ``พอยท์สะสมของคุณ'' แสดงยอดพอยท์ที่มีอยู่ ปุ่มสแกน QR Code รายการสถานที่ที่มี QR สำหรับรับพอยท์ และรายการของรางวัล ดังภาพที่ \ref{fig:points-member}

\screenshot{points-member}{หน้าพอยท์สะสมของสมาชิก}

\subsubsection{การสแกน QR Code เพื่อสะสมพอยท์ กรณีเป็นสมาชิก}

เมื่อไปถึงสถานที่ร่วมรายการ สมาชิกสแกนป้าย QR Code ด้วยปุ่มสแกนในหน้าพอยท์สะสม หรือใช้กล้องของโทรศัพท์สแกนโดยตรงก็ได้ ระบบจะขออนุญาตใช้กล้อง ดังภาพที่ \ref{fig:qr-scan}

\screenshot[0.7]{qr-scan}{หน้าต่างสแกน QR Code}

เมื่อสแกนสำเร็จ ระบบเพิ่มพอยท์เข้าบัญชีทันทีและแสดงผลการสแกน QR Code แต่ละจุดรับพอยท์ได้ครั้งเดียวต่อบัญชี หากสแกนซ้ำหรือ QR Code ไม่ถูกต้อง ระบบจะแสดงข้อความ ``สแกนไม่สำเร็จ'' พร้อมเหตุผล

\subsubsection{การแลกของรางวัล กรณีเป็นสมาชิก}

การแลกของรางวัลดำเนินการที่เคาน์เตอร์ โดยสมาชิกแจ้งชื่อหรือเบอร์โทรศัพท์ที่ใช้สมัครกับเจ้าหน้าที่ เจ้าหน้าที่จะหักพอยท์และมอบของรางวัลให้ผ่านหน้าผู้ดูแลระบบตามหัวข้อที่ \ref{sec:admin-guide}

\subsubsection{เมนูทริปของฉัน กรณีเป็นสมาชิก}

เมนู ``ทริปของฉัน'' แสดงแผนทริปที่บันทึกไว้ทั้งหมด ดังภาพที่ \ref{fig:my-trips} ผู้ใช้สามารถค้นหาตามชื่อทริป เปลี่ยนชื่อ เพิ่มในรายการโปรด ทำสำเนา และลบแผนได้ รวมถึงกด ``+ สร้างแผนใหม่'' เพื่อสร้างแผนเพิ่ม

\screenshot{my-trips}{หน้าทริปของฉัน}

\subsubsection{เมนูโปรไฟล์ กรณีเป็นสมาชิก}

ผู้ใช้เปิดหน้าโปรไฟล์ได้จากเมนูผู้ใช้มุมขวาบน ประกอบด้วยแท็บ 5 แท็บ ได้แก่ ภาพรวม ข้อมูล ประวัติ กิจกรรมของฉัน และความปลอดภัย ดังภาพที่ \ref{fig:profile-overview}

\screenshot{profile-overview}{หน้าโปรไฟล์ แท็บภาพรวม}

\begin{numlist}
\item ภาพรวม แสดงข้อมูลสรุปของผู้ใช้ ยอดพอยท์ และปุ่มลัด ``สแกน QR สะสมพอยท์'' และ ``ดูประวัติพอยท์''
\item ข้อมูล ใช้แก้ไขข้อมูลส่วนตัว เช่น คำนำหน้า ชื่อ นามสกุล วันเกิด อาชีพ เบอร์โทรศัพท์ และที่อยู่ แล้วกด ``บันทึกการเปลี่ยนแปลง'' ดังภาพที่ \ref{fig:profile-info}
\item ประวัติ แสดงประวัติพอยท์ทั้งการสแกน QR การแลกรางวัล และการปรับพอยท์ พร้อมตัวกรอง ดังภาพที่ \ref{fig:profile-history}
\item กิจกรรมของฉัน ใช้ส่งกิจกรรมใหม่เข้าระบบด้วยปุ่ม ``+ ส่งกิจกรรมใหม่'' และติดตามสถานะรออนุมัติ อนุมัติแล้ว หรือไม่อนุมัติ ดังภาพที่ \ref{fig:profile-events}
\item ความปลอดภัย ใช้เปลี่ยนรหัสผ่าน เปลี่ยนอีเมล และลบบัญชี ดังภาพที่ \ref{fig:profile-security}
\end{numlist}

\screenshot{profile-info}{หน้าโปรไฟล์ แท็บข้อมูล}

\screenshot{profile-history}{หน้าโปรไฟล์ แท็บประวัติ}

\screenshot{profile-events}{หน้าโปรไฟล์ แท็บกิจกรรมของฉัน}

\screenshot{profile-security}{หน้าโปรไฟล์ แท็บความปลอดภัย}

% -------------------------------------------------------------
\subsection{เข้าใช้งานกรณีเข้าสู่ระบบของผู้ดูแลระบบ}
\label{sec:admin-guide}
% -------------------------------------------------------------

\subsubsection{เมนูเข้าสู่ระบบ กรณีเป็นผู้ดูแลระบบ}

ผู้ดูแลระบบเข้าสู่ระบบที่ https://dinokhonkaen.app/admin/login หรือคลิกลิงก์ ``สำหรับผู้ดูแลระบบ'' ที่ส่วนท้ายของเว็บไซต์ ดังภาพที่ \ref{fig:admin-login} จากนั้นกรอกอีเมลและรหัสผ่านของบัญชีผู้ดูแลระบบแล้วกด ``เข้าสู่ระบบ'' บัญชีผู้ดูแลระบบสร้างโดยนักพัฒนาตามหัวข้อที่ \ref{sec:create-admin}

\screenshot[0.7]{admin-login}{หน้าเข้าสู่ระบบสำหรับผู้ดูแลระบบ}

\subsubsection{แถบเมนูและแดชบอร์ด}

หลังเข้าสู่ระบบ ระบบแสดงแถบเมนูด้านข้างที่จัดเป็นกลุ่ม ดังนี้ ภาพรวม (แดชบอร์ด) เนื้อหา (สถานที่ กิจกรรม ฐานความรู้แชทบอท) คิวตรวจสอบ (รายงานข้อมูล คำขอกิจกรรม) พอยท์และรางวัล (QR Code ของรางวัล แลกที่เคาน์เตอร์) และผู้ใช้และการใช้งาน (ผู้ใช้ ทริป ความพึงพอใจ) ซึ่งชื่อเมนูตรงกับที่แสดงในระบบ ส่วนแดชบอร์ดแบ่งเป็น 3 แท็บ ได้แก่ ภาพรวม ผู้ใช้และรางวัล และแผนทริป โดยแท็บภาพรวมแสดงจำนวนสถานที่ กิจกรรม ผู้ใช้ทั้งหมด ผู้สมัครใหม่ และการสแกน QR Code ใน 7 วัน รวมถึงจำนวนฐานความรู้ QR Code ของรางวัล และพอยท์ที่แจกไปแล้ว นอกจากนี้ยังมีส่วนอัปเดตดัชนีค้นหาของแชตบอตและเวอร์ชันของระบบ ดังภาพที่ \ref{fig:admin-dashboard}

\screenshot{admin-dashboard}{หน้าแดชบอร์ดผู้ดูแลระบบ}

\subsubsection{เมนูสถานที่}

ผู้ดูแลระบบจัดการข้อมูลสถานที่ได้จากเมนู ``สถานที่'' ดังภาพที่ \ref{fig:admin-places} โดยเพิ่มสถานที่ใหม่ แก้ไขข้อมูล ลบสถานที่ และนำเข้า/ซิงก์ข้อมูลจาก Google Places ได้ ฟอร์มเพิ่มและแก้ไขข้อมูลสถานที่แสดงดังภาพที่ \ref{fig:admin-place-form} ในการซิงก์ ผู้ดูแลระบบสามารถล็อกฟิลด์ที่แก้เองไว้ได้ ค่าใหม่จาก Google ของฟิลด์ที่ล็อกจะถูกพักไว้ให้ตรวจสอบก่อน และการซิงก์แต่ละสถานที่มีค่าใช้จ่ายตามการเรียก Google Places API

\screenshot{admin-places}{หน้าจัดการสถานที่}

\screenshot{admin-place-form}{หน้าฟอร์มเพิ่มและแก้ไขข้อมูลสถานที่}

\subsubsection{เมนูกิจกรรม}

เมนู ``กิจกรรม'' ใช้เพิ่ม แก้ไข และลบกิจกรรมและเทศกาล ดังภาพที่ \ref{fig:admin-events} ในฟอร์มกิจกรรมมีเครื่องมือช่วยดึงข้อมูลจากข้อความโพสต์ด้วย AI ให้ผู้ดูแลระบบวางข้อความโพสต์ แล้วระบบจะกรอกฟอร์มให้อัตโนมัติ ดังภาพที่ \ref{fig:admin-event-extract} ผู้ดูแลระบบควรตรวจสอบข้อมูลที่ AI กรอกให้ก่อนกดบันทึกทุกครั้ง

\screenshot{admin-events}{หน้าจัดการกิจกรรม}

\screenshot{admin-event-extract}{หน้าฟอร์มกิจกรรมพร้อมเครื่องมือดึงข้อมูลจากข้อความโพสต์ด้วย AI}

\subsubsection{เมนูฐานความรู้แชตบอต}

ผู้ดูแลระบบเพิ่ม แก้ไข ปักหมุด และเปิด/ปิดการใช้งานเอกสารความรู้ที่แชตบอตใช้อ้างอิง ดังภาพที่ \ref{fig:admin-knowledge} หลังเพิ่มหรือแก้ไขข้อมูล ต้องให้ระบบสร้าง Embedding ใหม่ เพื่อให้แชตบอตค้นพบข้อมูลล่าสุด โดยกดปุ่ม ``อัปเดตดัชนีค้นหา'' ในการ์ด ``ดัชนีค้นหาแชตบอท (Embedding)'' บนหน้าแดชบอร์ด

\screenshot{admin-knowledge}{หน้าจัดการฐานความรู้แชตบอต}

\subsubsection{เมนูรายงานข้อมูลและคำขอกิจกรรม}

เมนู ``รายงานข้อมูล'' แสดงคิวรายงานข้อมูลสถานที่และกิจกรรมที่ผู้ใช้แจ้งว่าไม่ถูกต้อง ผู้ดูแลระบบตรวจสอบและดำเนินการ ดังภาพที่ \ref{fig:admin-reports} ส่วนเมนู ``คำขอกิจกรรม'' แสดงกิจกรรมที่สมาชิกส่งเข้ามา ผู้ดูแลระบบเลือก ``อนุมัติ'' เพื่อแสดงกิจกรรมบนหน้าเว็บ หรือ ``ไม่อนุมัติ'' พร้อมระบุเหตุผลที่ผู้ขอจะเห็น ดังภาพที่ \ref{fig:admin-event-requests}

\screenshot{admin-reports}{หน้ารายงานข้อมูล}

\screenshot{admin-event-requests}{หน้าคำขอกิจกรรม}

\subsubsection{เมนู QR Code และของรางวัล}

เมนู ``QR Code'' ใช้สร้าง QR Code ผูกกับสถานที่และกำหนดจำนวนพอยท์ พร้อมดูตัวอย่างและพิมพ์ป้าย ดังภาพที่ \ref{fig:admin-qr} ส่วนเมนู ``ของรางวัล'' ใช้เพิ่มและแก้ไขรายการของรางวัล ราคาพอยท์ และจำนวนคงเหลือ ดังภาพที่ \ref{fig:admin-rewards}

\screenshot{admin-qr}{หน้าจัดการ QR Code}

\screenshot{admin-rewards}{หน้าจัดการของรางวัล}

\subsubsection{เมนูแลกที่เคาน์เตอร์}

เมื่อสมาชิกมาแลกของรางวัล เจ้าหน้าที่เปิดเมนู ``แลกที่เคาน์เตอร์'' ดังภาพที่ \ref{fig:admin-redeem} แล้วดำเนินการ 2 ขั้นตอน คือ เลือกผู้ใช้ และเลือกของรางวัล จากนั้นกดยืนยันเมื่อพร้อมมอบของ ระบบจะหักพอยท์ตามราคาของรางวัล หากบันทึกผิดสามารถยกเลิกรายการจากประวัติได้ ซึ่งระบบจะคืนพอยท์ให้ผู้ใช้และคืนของรางวัล 1 ชิ้นกลับเข้าจำนวนคงเหลือ

\screenshot{admin-redeem}{หน้าแลกของรางวัลที่เคาน์เตอร์}

\subsubsection{เมนูผู้ใช้ ทริป และความพึงพอใจ}

เมนู ``ผู้ใช้'' ใช้ดูรายละเอียดและจัดการบัญชีผู้ใช้ ดังภาพที่ \ref{fig:admin-users} เมนู ``ทริป'' ใช้ค้นหาและดูแผนทริปที่ผู้ใช้สร้างไว้ ดังภาพที่ \ref{fig:admin-trips} และเมนู ``ความพึงพอใจ'' แสดงผลแบบสอบถามความพึงพอใจ ได้แก่ คะแนนเฉลี่ย จำนวนผู้ตอบ ข้อเสนอแนะ และคำตอบรายคน ดังภาพที่ \ref{fig:admin-survey}

\screenshot{admin-users}{หน้าจัดการผู้ใช้}

\screenshot{admin-trips}{หน้าดูทริปของผู้ใช้}

\screenshot{admin-survey}{หน้าผลความพึงพอใจ}


\chapter{การติดตั้งระบบสำหรับพัฒนา}
% =============================================================
% บทที่ 3 การติดตั้งระบบสำหรับพัฒนา
% ข้อมูลตรวจจาก main repo: README.md, CLAUDE.md, .env.example ของทั้ง 3 ส่วน, backend/scripts
% =============================================================

\section{ความต้องการของระบบ}

\subsection{ระบบปฏิบัติการ}

คู่มือนี้อธิบายการติดตั้งบนระบบปฏิบัติการ Windows 11

\subsection{ซอฟต์แวร์และบริการที่ต้องเตรียม}

\begin{parenlist}
\item Node.js เวอร์ชัน 20 ขึ้นไป (ระบบ CI ใช้เวอร์ชัน 20) สำหรับส่วนติดต่อผู้ใช้งานและบริการข้อมูลหลัก
\item Python เวอร์ชัน 3.11 สำหรับ AI Service
\item Git สำหรับดาวน์โหลดซอร์สโค้ดจาก GitHub
\item Visual Studio Code สำหรับเปิดและแก้ไขโครงงาน
\item บัญชี Supabase สำหรับฐานข้อมูล PostgreSQL
\item บัญชี Cloudinary สำหรับเก็บรูปภาพ
\item บัญชี Gmail ที่เปิดการยืนยันตัวตน 2 ขั้นตอนและสร้าง App Password สำหรับส่งอีเมล
\item โปรเจกต์ Google Cloud ที่เปิดใช้ Places API (New) และ Maps JavaScript API
\item คีย์ KKU AI Gateway สำหรับเรียกโมเดลภาษาขนาดใหญ่
\item คีย์ Serper API (ไม่บังคับ) สำหรับสคริปต์สร้างคำอธิบายสถานที่
\end{parenlist}

\section{ขั้นตอนการติดตั้งระบบสำหรับนักพัฒนา}

\subsection{การติดตั้ง Node.js}

\begin{numlist}
\item เข้าเว็บไซต์ https://nodejs.org/ แล้วเลือกดาวน์โหลดเวอร์ชัน LTS ดังภาพที่ \ref{fig:install-nodejs}
\item เปิดไฟล์ติดตั้งที่ดาวน์โหลดมา แล้วกด Next จนติดตั้งเสร็จสิ้น
\item เปิด PowerShell แล้วตรวจสอบการติดตั้งด้วยคำสั่งต่อไปนี้ ซึ่งควรแสดงหมายเลขเวอร์ชัน
\end{numlist}

\screenshot{install-nodejs}{หน้าดาวน์โหลด Node.js}

\begin{lstlisting}[numbers=none]
node -v
npm -v
\end{lstlisting}

\subsection{การติดตั้ง Python}

\begin{numlist}
\item เข้าเว็บไซต์ https://www.python.org/downloads/windows/ แล้วดาวน์โหลดตัวติดตั้ง Windows installer (64-bit) ของรุ่น 3.11 ขึ้นไป ดังภาพที่ \ref{fig:install-python}
\item ขณะติดตั้งให้เลือกตัวเลือก ``Add python.exe to PATH'' แล้วกด Install Now
\item ตรวจสอบการติดตั้งด้วยคำสั่ง \texttt{python --version}
\end{numlist}

\screenshot{install-python}{หน้าดาวน์โหลด Python}

\subsection{การติดตั้ง Git และ Visual Studio Code}

\begin{numlist}
\item ดาวน์โหลดและติดตั้ง Git จาก https://git-scm.com/downloads
\item ดาวน์โหลดและติดตั้ง Visual Studio Code จาก https://code.visualstudio.com/
\end{numlist}

\subsection{การดาวน์โหลดซอร์สโค้ดจาก GitHub}

เปิด PowerShell ไปยังโฟลเดอร์ที่ต้องการเก็บโครงงาน แล้วโคลน repository

\begin{lstlisting}[numbers=none]
git clone https://github.com/Poeir/dino-trip.git
cd dino-trip
\end{lstlisting}

โครงสร้างโฟลเดอร์ที่ได้ประกอบด้วย 3 ส่วนที่รันแยกจากกัน ดังตารางที่ \ref{tab:folders}

\begin{table}[H]
\caption{โครงสร้างโฟลเดอร์ของโครงงาน}
\label{tab:folders}
\centering
\begin{tabular}{|l|l|p{5cm}|}
\hline
\textbf{โฟลเดอร์} & \textbf{เทคโนโลยี} & \textbf{หน้าที่} \\
\hline
\texttt{frontend} & React 18 + Vite & หน้าเว็บที่ผู้ใช้และผู้ดูแลระบบใช้งาน \\
\hline
\texttt{backend} & Node.js + Express & REST API ฐานข้อมูล และสคริปต์นำเข้าข้อมูล \\
\hline
\texttt{chatbot-service} & Python + FastAPI & แชตบอต ระบบวางแผนทริป และการสร้าง Embedding \\
\hline
\end{tabular}
\end{table}

\subsection{การสร้างฐานข้อมูลบน Supabase}

\begin{numlist}
\item เข้าเว็บไซต์ https://supabase.com/ ลงชื่อเข้าใช้ แล้วสร้างโปรเจกต์ใหม่ พร้อมบันทึกรหัสผ่านฐานข้อมูลที่กำหนดไว้
\item เปิดหน้า Project Settings > Database > Connection string แล้วคัดลอก Session pooler (พอร์ต 5432) ไว้ใช้เป็นค่า \texttt{DATABASE\_URL} ไม่ควรใช้ Transaction mode (พอร์ต 6543) เพราะ backend เป็นเซิร์ฟเวอร์ที่รันต่อเนื่อง
\item เปิดหน้า Project Settings > API แล้วคัดลอก Project URL และ service\_role key ไว้ใช้กับ AI Service
\item สร้างโครงสร้างตารางด้วยการนำไฟล์ migration ไปใช้กับฐานข้อมูลใหม่ ดังคำสั่งในหัวข้อการนำเข้าข้อมูลตั้งต้น
\end{numlist}

\subsection{การติดตั้งไลบรารีของแต่ละส่วน}

เปิดโฟลเดอร์โครงงานด้วย Visual Studio Code (\texttt{code .}) แล้วใช้ Terminal ติดตั้งไลบรารีของแต่ละส่วนตามคำสั่งต่อไปนี้

\begin{lstlisting}[numbers=none]
cd backend
npm install

cd ../frontend
npm install

cd ../chatbot-service
python -m venv venv
venv\Scripts\activate
pip install -r requirements.txt
\end{lstlisting}

\subsection{การตั้งค่าไฟล์ .env}

แต่ละส่วนมีไฟล์ \texttt{.env.example} เป็นต้นแบบ ให้คัดลอกเป็นไฟล์ \texttt{.env} ในโฟลเดอร์เดียวกัน แล้วกรอกค่าตามตารางที่ \ref{tab:env-vars} ไฟล์ \texttt{.env} มีข้อมูลลับ จึงห้ามนำขึ้น Git

\begin{lstlisting}[numbers=none]
cd backend
copy .env.example .env
cd ../frontend
copy .env.example .env
cd ../chatbot-service
copy .env.example .env
\end{lstlisting}

\begin{table}[H]
\caption{ตัวแปรสภาพแวดล้อมที่ต้องกำหนดในไฟล์ .env}
\label{tab:env-vars}
\centering
\begin{tabular}{|l|l|p{5.2cm}|}
\hline
\textbf{ส่วนของระบบ} & \textbf{ตัวแปร} & \textbf{คำอธิบาย} \\
\hline
backend & \texttt{DATABASE\_URL} & connection string ของ Postgres (Session pooler พอร์ต 5432) \\
\hline
backend & \texttt{FRONTEND\_ORIGIN} & origin ของ frontend ต้องตรงกันทุกตัวอักษรและไม่มี / ท้าย เช่น \texttt{http://localhost:5173} \\
\hline
backend & \texttt{CHATBOT\_SERVICE\_URL} & ที่อยู่ของ AI Service เช่น \texttt{http://localhost:8000} \\
\hline
backend & \texttt{CLOUDINARY\_CLOUD\_NAME}, \texttt{CLOUDINARY\_API\_KEY}, \texttt{CLOUDINARY\_API\_SECRET} & ค่าจาก Cloudinary Dashboard > Settings > API Keys \\
\hline
backend & \texttt{GMAIL\_USER}, \texttt{GMAIL\_APP\_PASSWORD} & บัญชี Gmail และ App Password สำหรับส่งอีเมลยืนยัน \\
\hline
backend & \texttt{GOOGLE\_PLACES\_API\_KEY} & คีย์สำหรับนำเข้าและซิงก์สถานที่ (ต้องเปิด Places API (New) และผูก billing) \\
\hline
backend & \texttt{PORT} & พอร์ตของ backend ค่าเริ่มต้น 4000 \\
\hline
frontend & \texttt{VITE\_API\_URL} & ที่อยู่ของ backend เช่น \texttt{http://localhost:4000} \\
\hline
frontend & \texttt{VITE\_CHATBOT\_SERVICE\_URL} & ที่อยู่ของ AI Service เช่น \texttt{http://localhost:8000} \\
\hline
frontend & \texttt{VITE\_GOOGLE\_MAPS\_API\_KEY} & คีย์ Maps JavaScript API ควรจำกัดด้วย HTTP referrer \\
\hline
chatbot-service & \texttt{KKU\_API\_KEY} & คีย์ของ KKU AI Gateway \\
\hline
chatbot-service & \texttt{SUPABASE\_URL}, \texttt{SUPABASE\_SERVICE\_ROLE\_KEY} & ใช้ค้นข้อมูลสถานที่และค้นหาเชิงเวกเตอร์ \\
\hline
chatbot-service & \texttt{MODEL\_NAME}, \texttt{TRIP\_PLANNER\_MODEL\_NAME} & โมเดลของแชตบอตและระบบวางแผนทริป (ต้องตั้งบน production) \\
\hline
chatbot-service & \texttt{CORS\_ORIGINS} & origin ที่เรียก AI Service ได้ คั่นด้วย , และไม่มี / ท้าย \\
\hline
\end{tabular}
\end{table}

\subsection{การนำเข้าข้อมูลตั้งต้น}

ขั้นตอนนี้ทำเพียงครั้งเดียวตอนติดตั้งฐานข้อมูลใหม่ โดยใช้ Terminal ที่โฟลเดอร์ \texttt{backend}

\begin{lstlisting}[numbers=none]
cd backend
npx supabase db push --yes
npm run fetch:places
npm run import:places
npm run seed:events-knowledge
\end{lstlisting}

ความหมายของแต่ละคำสั่งมีดังนี้

\begin{numlist}
\item \texttt{npx supabase db push --yes} สร้างตารางและฟังก์ชันตามไฟล์ migration ใช้กับฐานข้อมูลใหม่เท่านั้น เพราะฐานข้อมูลที่ใช้งานจริงอยู่แล้วไม่มีประวัติ migration ทำให้คำสั่งนี้รัน migration เก่าซ้ำทั้งหมด
\item \texttt{npm run fetch:places} ดึงข้อมูลสถานที่จาก Google Places API ลงไฟล์ \texttt{backend/data/places.json} ซึ่งมีค่าใช้จ่ายตามจำนวนสถานที่
\item \texttt{npm run import:places} นำเข้าข้อมูลจากไฟล์ดังกล่าวเข้าฐานข้อมูล
\item \texttt{npm run seed:events-knowledge} เพิ่มข้อมูลกิจกรรมและฐานความรู้ตั้งต้น
\end{numlist}

หลังนำเข้าข้อมูลใหม่ทุกครั้ง ต้องสร้าง Embedding ใหม่ เพื่อให้แชตบอตและระบบวางแผนทริปค้นพบข้อมูล

\begin{lstlisting}[numbers=none]
cd chatbot-service
venv\Scripts\activate
python scripts/embed_content.py
\end{lstlisting}

\subsection{การสร้างบัญชีผู้ดูแลระบบ}
\label{sec:create-admin}

บัญชีผู้ดูแลระบบสร้างด้วยสคริปต์ที่โฟลเดอร์ \texttt{backend} โดยระบุอีเมลและรหัสผ่านที่ต้องการ

\begin{lstlisting}[numbers=none]
cd backend
node scripts/create-admin.js <email> <password>
\end{lstlisting}

เมื่อสำเร็จจะแสดงข้อความ ``Admin account ready'' และสามารถเข้าสู่ระบบที่หน้า \texttt{/admin/login} ได้

\subsection{การรันระบบ}

เปิด Terminal ใน Visual Studio Code แยก 3 หน้าต่าง แล้วรันแต่ละส่วนตามคำสั่งต่อไปนี้

\begin{lstlisting}[numbers=none]
# Terminal 1: Backend (http://localhost:4000)
cd backend
npm run dev

# Terminal 2: AI Service (http://localhost:8000)
cd chatbot-service
venv\Scripts\activate
uvicorn src.main:app --reload --port 8000

# Terminal 3: Frontend (http://localhost:5173)
cd frontend
npm run dev
\end{lstlisting}

เมื่อทุกส่วนทำงานแล้ว ให้เปิดเบราว์เซอร์ไปที่ \texttt{http://localhost:5173} จะปรากฏหน้าแรกของไดโน และสามารถดูเอกสาร API ของ AI Service ได้ที่ \texttt{http://localhost:8000/docs}

\subsection{การนำระบบขึ้นใช้งานจริง}

ระบบที่ใช้งานจริงที่ https://dinokhonkaen.app รันบนเซิร์ฟเวอร์เดียวด้วย Docker Compose ในโฟลเดอร์ \texttt{deploy} โดยมี Caddy ให้บริการ HTTPS เสิร์ฟหน้าเว็บ และส่งต่อเส้นทาง \texttt{/api} ไปยัง backend และ \texttt{/chat}, \texttt{/trip/llm} ไปยัง AI Service ทั้งนี้ การรวมโค้ดเข้ากิ่ง \texttt{main} บน GitHub จะทำให้ GitHub Actions เชื่อมต่อเซิร์ฟเวอร์และอัปเดตระบบอัตโนมัติ


\end{document}
